\documentclass[11pt]{article}
\usepackage[a4paper,margin=28mm]{geometry}
\usepackage{amsmath,amssymb,amsthm,xcolor,booktabs,array}
\newcommand{\PROVED}{\textcolor{green!40!black}{\textbf{PROVED}}}
\newcommand{\OPEN}{\textcolor{red!70!black}{\textbf{OPEN}}}
\newcommand{\AUDIT}{\textcolor{orange!70!black}{\textbf{AUDIT}}}
\newtheorem{proposition}{Proposition}
\newtheorem{lemma}{Lemma}
\title{Row (d): audit of the possible hidden closure (v95)}
\date{13 August 2026}

\begin{document}
\maketitle

\section*{Purpose}
This note audits whether rows (a)--(c), the mixed-line construction in
\texttt{main.tex}, or the identities collected through v94 already contain a
non-circular proof of row (d).  Statements are separated strictly into
\PROVED, \OPEN, and \AUDIT.

\section{The one-line closure that would work}

Let $\mathbb L_{\mathrm{comp}}^{\bar0}=\mathcal H_+^0$ and
$\mathbb L_{\mathrm{comp}}^{\bar1}=\mathcal H_-^0$ be the completed
super-representation of row (c), and let $\pi_-$ denote the action on the odd
quotient.  The even quotient consists of the two Tate characters.

\begin{proposition}[Hilbert--Rosati closure criterion --- \PROVED]
Suppose that the odd quotient admits a Hilbert completion for which
\[
 \pi_-(f^\vee)=\pi_-(f)^*
\]
on the test algebra and $\pi_-(f)\pi_-(f)^*$ is trace class.  Then, for every
primitive $f$,
\[
 B_{\rm nuc}(f,f)=-\operatorname{Tr}
 \bigl(\pi_-(f)\pi_-(f)^*\bigr)\leq0.
\]
Consequently row (d), condition $G$, and Weil positivity follow.
\end{proposition}

\begin{proof}
Row (c) expresses $B_{\rm nuc}(f,f)$ as the supertrace of the action of
$f\star f^\vee$.  On the even quotient this action is evaluation at the two
Tate characters, hence its trace is
$|\widehat f(0)|^2+|\widehat f(1)|^2$.  It vanishes for primitive $f$.
Under the asserted adjoint identity the odd action is
$\pi_-(f)\pi_-(f)^*$, a positive trace-class operator.  The odd parity enters
the supertrace with a minus sign, which proves the formula and the inequality.
\end{proof}

This is the genuine potential hidden closure.  It uses all three previous
rows: the test algebra and geometry from (a), the correspondence action from
(b), and the supertrace comparison from (c).

\section{Why the hypothesis is not already in rows (a)--(c)}

\begin{proposition}[Type audit --- \PROVED]
Rows (a)--(c) prove an algebraic Tate involution and Hermitian symmetry of
$B_{\rm nuc}$, but they do not construct a positive Hilbert norm on
$\mathcal H_-^0$ for which the odd action is a $*$-representation.
\end{proposition}

\begin{proof}
Row (a) constructs Hilbert/cotangent structures on the external ruling
sector.  Its added nuclear direction is explicitly not asserted to be a
mixed Cartier divisor.  Row (b) constructs the monoidal action
$\Gamma_m\Gamma_n\simeq\Gamma_{mn}$ and an adjoint Witt dynamics, but no
comparison identifies that adjoint with the Tate involution on the completed
odd Poisson quotient.  Row (c) places $\mathcal H_-^0$ in the category of
nuclear Fr\'echet scaling representations and proves summability and a
supertrace formula.  A nuclear Fr\'echet representation is not, from these
facts alone, a Hilbert $*$-representation.  Thus the adjoint identity used in
the preceding proposition is absent.
\end{proof}

\paragraph{Status.} The construction of such a Hilbert--Rosati completion,
from the source geometry and before using the spectral sign, is \OPEN.  If
one defines its positive cone or norm using the sign of $B_{\rm nuc}$, the
argument is circular and is classified \AUDIT.

\section{Audit of the object called $D_f$ in the mixed-line proposal}

The proposal defines $\lambda_{\rm fin}(f,g)$ from the convolution
$f\star g^\vee$ and, for $g=f$, twists it by an operator $e^{-\Sigma_f}$ so
that the logarithmic covolume equals $B_{\rm nuc}(f,f)$.

\begin{proposition}[Quadratic determinant versus divisor class --- \PROVED]
The construction proves a determinant/covolume identity for the quadratic
pair $(f,f)$.  It does not, as presently written, construct an additive mixed
divisor class $f\mapsto D_f$ with a tensor law
\[
 D_{f+g}=D_f+D_g,
 \qquad \mathcal O(D+E)\simeq\mathcal O(D)\otimes\mathcal O(E),
\]
nor a section functor on its multiples.
\end{proposition}

\begin{proof}
The exponent defining $\lambda_{\rm fin}(f,f)$ contains
$(f\star f^\vee)(\pm\log n)$ and is therefore quadratic in $f$.
Replacing $f$ by $f+g$ produces the two cross terms
$f\star g^\vee$ and $g\star f^\vee$; hence the resulting normed line is not
the tensor product of the lines assigned separately to $f$ and $g$.
Moreover $e^{-\Sigma_f}$ is an operator metric whose trace records the
archimedean determinant; a metric on a rank-one line is scalar.  Passing to
its Fredholm determinant retains the scalar covolume but not an object of the
mixed Picard groupoid with sections.  Therefore the displayed covolume
identity constructs the desired self-intersection number, not the underlying
linear divisor class needed by Riemann--Roch.
\end{proof}

It follows that the notation $D_f$ in that paragraph must be read as a
quadratic determinant object unless an additional linearization theorem is
proved.  Treating it as the already constructed mixed divisor is \AUDIT.

\section{Audit of the theta and Sonin continuation}

\begin{proposition}[Exact status --- \PROVED]
The theta construction proves the four numerical Riemann--Roch/effectivity
axioms on the interior of the external ruling cone.  It does not prove them
on primitive mixed classes.  The Sonin Gram $M(f)$ is positive semidefinite,
but this does not imply the lower-frame estimate
\[
 \lambda_{\min}(M(f))>
 \frac{\operatorname{tr}M(f)+B_{\rm nuc}(f,f)}{|S_T|}.
\]
\end{proposition}

\begin{proof}
On primitive classes both ruling degrees vanish, so the ruling code has
bounded rank while mixed Riemann--Roch requires quadratic growth.  For the
Sonin metric, Gram positivity gives only $\lambda_{\min}(M(f))\geq0$.
The required right-hand side has no proved nonpositive sign, and positive
semidefiniteness supplies no comparison between the smallest eigenvalue and
that scalar.  Thus the displayed strict inequality is an additional theorem,
not a consequence of Gram positivity.
\end{proof}

The general-window lower-frame estimate is \OPEN.  Replacing it by positivity
of $M(f)$ is \AUDIT.

\section{Audit of v91--v94}

The Gamma--Fourier--Tate identities, first-loss reductions, Ward equations,
and causal formulations in v91--v94 establish exact equivalences and rule out
several false bridges.  None constructs either

\begin{enumerate}
\item the Hilbert--Rosati $*$-completion of $\mathcal H_-^0$; or
\item an additive mixed divisor with a section/effectivity functor satisfying
mixed Riemann--Roch.
\end{enumerate}

In particular, equations of the form (565), (661), and the Ward terminal
identity are analytic forms of the no-first-contact assertion.  They are not
independent inputs from rows (a)--(c).  Their equivalence to the terminal
positivity is \PROVED; their source-geometric realization remains \OPEN.

\section{Correct terminal ledger}

\begin{center}
\begin{tabular}{>{\raggedright\arraybackslash}p{.35\linewidth}
                >{\raggedright\arraybackslash}p{.16\linewidth}
                >{\raggedright\arraybackslash}p{.39\linewidth}}
\toprule
Item & Status & Reason \\
\midrule
Rows (a), (b), (c), under their stated typed contracts & \PROVED & Their
objects, actions, contacts, and nuclear trace comparison are constructed.\\
Ruling-sector RR/effectivity and Hodge signature & \PROVED & Complete on the
external two-degree sector.\\
Full Hilbert-space row (d) for $0<T\leq\log2$ & \PROVED & Certified local
operator theorem.\\
Quadratic covolume identity for the proposed mixed determinant & \PROVED & It
equals $B_{\rm nuc}(f,f)$ without asserting its sign.\\
Additive mixed divisor and its section functor & \OPEN & No linear tensor law,
effectivity theory, or mixed RR has been constructed.\\
Odd Hilbert--Rosati $*$-completion & \OPEN & This would close row (d) by the
one-line supertrace argument above.\\
General Sonin lower-frame bound & \OPEN & Gram positivity alone is
insufficient.\\
Global row (d), condition $G$ & \OPEN & Follows from either of the two
preceding genuine constructions, neither of which is proved.\\
\bottomrule
\end{tabular}
\end{center}

\section*{Single next theorem}

The most economical remaining target is:

\medskip
\noindent\textbf{Nuclear Hodge--Rosati theorem (\OPEN).}
Construct, directly from the Poisson quotient and the row-(a)/(b) source
geometry, a Hilbert completion of $\mathcal H_-^0$ such that
$\pi_-(f^\vee)=\pi_-(f)^*$ and the square is trace class, without defining
the norm from $B_{\rm nuc}$ or from the zeros.  The first proposition then
proves row (d) with no further inequality.

\end{document}
